% Generated draft: requires mathematical and source review.
\documentclass[11pt]{article}
\usepackage[margin=25mm]{geometry}
\usepackage{fontspec}
\setmainfont{Latin Modern Roman}
\usepackage{amsmath,amssymb,mathtools,unicode-math}
\setmathfont{Latin Modern Math}
\usepackage{enumitem,needspace,xurl,hyperref}
\setlength{\parindent}{0pt}
\setlength{\parskip}{5pt}
\setlength{\emergencystretch}{3em}
\newcommand{\N}{\mathbb{N}}
\newcommand{\Z}{\mathbb{Z}}
\newcommand{\Q}{\mathbb{Q}}
\newcommand{\R}{\mathbb{R}}
\newcommand{\C}{\mathbb{C}}
\newcommand{\F}{\mathbb{F}}
\newcommand{\A}{\mathbb{A}}
\newcommand{\PP}{\mathbb{P}}
\newcommand{\E}{\mathbb{E}}
\newcommand{\eps}{\varepsilon}
\newcommand{\dd}{\,\mathrm{d}}
\begin{document}
\section*{UnsolvedMath ID: 2973}
\textbf{Source problem code:} catalogue entry
\subsection{Definitions and mathematical statement}
Let $(X,J)$ be a closed almost complex $4$-manifold with its induced orientation. Write
\[c_1^2(X)=\langle c_1(TX,J)^2,[X]\rangle,\quad c_2(X)=\langle c_2(TX,J),[X]\rangle=e(X),\quad c_1^2(X)=2e(X)+3\sigma(X).\]
When $c_2(X)\ne0$, the source's Chern slope is $\lambda(X)=c_1^2(X)/c_2(X)=2+3\sigma(X)/e(X)$. Let $\Sigma_g$ be the closed oriented surface of genus $g$. The question asks for four separate suprema of $\lambda(X)$ as $X$ ranges over:
\begin{itemize}
\item (a) Symplectic smooth $\Sigma_g$-bundles over $\Sigma_h$ with $g,h\geq2$.
\item (b) Holomorphic smooth $\Sigma_g$-bundles over $\Sigma_h$ with $g,h\geq2$, meaning the bundle projection is holomorphic for suitable complex structures; no isotriviality requirement is added.
\item (c) Symplectic Lefschetz fibrations over $S^2$.
\item (d) Holomorphic Lefschetz fibrations over $S^2$.
\end{itemize}
In both (c) and (d), the critical set must be nonempty. A Lefschetz fibration has isolated standard positive nodal critical points; the source does not specify further minimality conventions or the allowed fibre genera in these two cases. The supremum need not be attained, and equality of the four answers is not assumed. For (a) and (b), $e(X)=4(g-1)(h-1)>0$. For the Lefschetz cases, the ratio still requires a nonzero denominator; the record does not explicitly give a positivity/exclusion convention for $c_2$, so no undefined slope or unstated minimal/general-type hypothesis is silently included.

The notes report that such bundles with $g\geq2$ and Lefschetz fibrations with critical points over $S^2$ admit symplectic structures, although not necessarily holomorphic ones [Joh86, Bay12b]. They cite [Hil00, Kot99] for the relation between existence of a complex structure on the total space and holomorphic realization of a surface bundle. They report signature zero for bundles if $g\leq2$ or $h\leq1$ [Mey73], and positive-signature symplectic examples for all but 19 candidate pairs $g\geq3,h\geq2$ [BK24b]. In particular, the source says a positive-signature $(g,h)=(3,2)$ example would violate the BMY slope bound, while a $(4,2)$ example with smallest possible positive signature would lie on the BMY line and be noncomplex. These are implications of hypothetical examples, not constructions asserted here; the source refers to [BK24b, Question 2] and Problem 4.90.

The source encodes the complex-surface Bogomolov--Miyaoka--Yau bound as $\lambda\leq3$ in its intended slope setting, and says bundles in (b) satisfy the strict inequality $\lambda<3$ because complex ball quotients cannot smoothly fibre over a surface. Strict inequality for every member alone does not determine its supremum. An unpublished Hamenstädt preprint announces $\lambda\leq3$ also for (a); it is retained as an announcement, not upgraded to a verified universal bound. The notes' largest known slopes for (a) and (b) are $8/3$ from Catanese--Rollenske. They cite Bryan--Donagi for holomorphic positive-signature bundles with $h=2$ and large $g$, and Kazuhiro Konno's Japanese-language works for $g=3$ with unspecified, presumably large base genus, without supplying a precise Konno reference.

For (c),(d), the source reports recent positive-signature symplectic Lefschetz examples over $S^2$ [BH24c]. It compares the Cartwright--Steger holomorphic Lefschetz fibration with slope $3$ over $T^2$ [KY21], whose different base excludes it from (c),(d). The supplied triage checked 2026-08-17 labels the four-way supremum programme partially solved / PARTIAL-PROGRESS, reports construction families and known bounds but no complete computation, and the registry is open.
\subsection{Short English statement}
Find the largest possible Chern slope separately for symplectic and holomorphic surface bundles, and for symplectic and holomorphic Lefschetz fibrations over the sphere. The source reports lower examples and qualified upper-bound claims, while the four sharp suprema remain undetermined and some slope-domain conventions are missing.
\subsection{Sources}
\begin{itemize}
\item Local UnsolvedMath registry, permanent ID 2973, source code catalogue entry, ``Kirby Problem 4.97''; supplied four categories, geography notes, and bibliography. Source Problem 4.90 is a cross-reference, not a permanent-ID identification.
\item \emph{K3: A New Problem List in Low-Dimensional Topology}, Problem 4.97 (2026): \url{https://aimath.org/pastworkshops/kirbylistrep.pdf}.
\item [Joh86] F. E. A. Johnson, ``A class of non-Kählerian manifolds,'' \emph{Mathematical Proceedings of the Cambridge Philosophical Society} 100(3) (1986), 519--521, doi:10.1017/S030500410006624X.
\item [Bay12b] R. İnanç Baykur, ``Non-holomorphic surface bundles and Lefschetz fibrations,'' \emph{Mathematical Research Letters} 19(3) (2012), 567--574, doi:10.4310/MRL.2012.v19.n3.a5.
\item [Hil00] Jonathan A. Hillman, ``Complex surfaces which are fibre bundles,'' \emph{Topology and Its Applications} 100(2--3) (2000), 187--191, doi:10.1016/S0166-8641(98)00085-6.
\item [Kot99] D. Kotschick, ``On regularly fibered complex surfaces,'' in \emph{Proceedings of the Kirbyfest (Berkeley, California, 1998)}, Geometry and Topology Monographs 2 (1999), 291--298, Geometry and Topology Publications, Coventry, doi:10.2140/gtm.1999.2.291.
\item [Mey73] Werner Meyer, ``Die Signatur von Flächenbündeln,'' \emph{Mathematische Annalen} 201 (1973), 239--264, doi:10.1007/BF01427946.
\item [BK24b] R. İnanç Baykur and Mustafa Korkmaz, ``Geography of surface bundles over surfaces,'' \emph{Mathematische Annalen} 390(4) (2024), 5793--5817, doi:10.1007/s00208-024-02899-5; background also cites Question 2.
\item [Ham20] Ursula Hamenstädt, ``Signature of surface bundles and bounded cohomology'' (2020), arXiv:2011.05792. The source expressly calls this an unpublished preprint announcing the bound.
\item [CR09] Fabrizio Catanese and Sönke Rollenske, ``Double Kodaira fibrations,'' \emph{Journal für die reine und angewandte Mathematik} 628 (2009), 205--233, doi:10.1515/CRELLE.2009.024.
\item [BD02] Jim Bryan and Ron Donagi, ``Surface bundles over surfaces of small genus,'' \emph{Geometry and Topology} 6 (2002), 59--67, doi:10.2140/gt.2002.6.59.
\item [BH24c] R. İnanç Baykur and Noriyuki Hamada, ``Lefschetz fibrations with arbitrary signature,'' \emph{Journal of the European Mathematical Society} 26(8) (2024), 2837--2895, doi:10.4171/jems/1326.
\item [KY21] Vincent Koziarz and Sai-Kee Yeung, ``Stability of the Albanese fibration on the Cartwright-Steger surface,'' \emph{Taiwanese Journal of Mathematics} 25(2) (2021), 251--256, doi:10.11650/tjm/201108.
\item The background attributes a $g=3$ holomorphic positive-signature construction to Kazuhiro Konno, in Japanese-language works, but supplies no full reference or base genus; none is invented.
\item Supplied literature triage checked 2026-08-17, PARTIAL-PROGRESS; it reports no complete four-way supremum calculation, preserves the damaged lambda display in the source, and cautions that dated triage does not exclude later or unindexed results.
\end{itemize}
\end{document}
